\section{Threats to Validity}\label{sec:threats}

\emph{Internal.} On CIC-IoT2023 the same author identified the attackers, wrote the labelling rule and made all design choices. The labels were revised once after error analysis (the sixth DDoS source), towards the model's predictions; a later blind audit in both directions found one more error, whose correction also favours the model (Section~\ref{sec:labels}). The audit's rater is a program, not a second person, and four of its rules were fixed after its first run (one a post-hoc criterion); we release its blind sample for a manual check. Our CIC-IoT2023 results are scored against our own labels; IoT-23 supports the cost of file-name labels with the authors' labels but does not validate ours, and its file-name labels depend on our family mapping.

\emph{Construct.} Macro-F1 weighs tiny classes (Mirai, MITM-ARP) like large ones and, on IoT-23, missing classes; we report per-class results with host-pair counts. Without MITM-ARP, which flow features cannot see by the label's definition, macro-F1 is 0.826 instead of 0.725. False blocks per hour extrapolate from ten minutes of live traffic; with persistent counters the reset is shown to hide some (Section~\ref{sec:response}).

\emph{External.} Labels and splits were audited on two datasets, but the live and response steps only on CIC-IoT2023, with one capture session per attack class, so they measure generalisation within a session, not across networks. We do not test across datasets, where results are known to be asymmetric~\cite{ali2026siamxbert}. Large CIC-IoT2023 captures were trimmed, so a temporal split covers only their start. Live results hold only at the achieved replay rates, on one laptop under WSL2, and for floods that did not fill the capture buffer. Replay is open-loop, so ``packets dropped'' is counterfactual: drops cause no retransmissions, and attackers do not adapt.

\emph{Statistical conclusion.} The enforced settings ran once, and each pinned CPU budget three times. The SSH-to-RTSP change rests on three host pairs, and the CIC-IoT2023 gap beyond it is small and not established at capture level. LOCO tests one held-out capture at a time, and its ranges span ten captures per dataset. ``No packets were lost'' refers only to these runs.

\section{Conclusion}\label{sec:conclusion}

Each step of the audit exposed a problem that a standard evaluation hid, from labels to blocking, and once in our own live results, where flow splitting left macro-F1 unchanged. The four checks that caught them are the reusable result. Future work includes a blocking criterion that weighs attack volume, not only flow verdicts, longer live runs, a second live dataset, early classification~\cite{pekar2025early} and edge hardware.

\subsection{Artifact Availability}

The repository \url{https://github.com/muhammadusamahoyrr/EdgeFlow-IDS} (MIT licence) holds the code, attacker evidence, configuration, result files and audit scripts. The scripts recompute the offline numbers from the datasets; the live numbers also need the live flow logs, which are not in the repository. All random choices (split, trees, NSGA-II, bootstrap) use seed 0; live runs are timing-dependent. The NFStream fix is included (\nolinkurl{src/patch_nfstream.py}, diff in \nolinkurl{patches/}).
% PLACEHOLDER for the author (do not compile in until minted): add one sentence naming the tagged
% release used for this paper and its Zenodo DOI, e.g. "The exact version used here is release
% <tag>, archived at Zenodo (DOI <to be minted>)."
The CIC-IoT2023 and IoT-23 data are available from their authors~\cite{neto2023ciciot2023,garcia2020iot23}.
